\section{结论}

本文围绕四个子问题建立了统一的“水流事件场--时间依赖路权--动态重规划”框架，主要结论如下。

（1）单源条件下，矿井 1 的 663 个端点、977 条巷道全部进入水流影响范围；矿井 2 有 28 个端点可达、39 条巷道充满。累计体积法避免了分叉流量被重复用于传播和蓄水的问题。

（2）首次通知后，六名矿工均可到达出口 1，到达时刻为 $12.17$--$19.20\,\mathrm{min}$。穿越途中若水态改变，算法会即时切换速度并检查 $0.3\,\mathrm{m}$ 阈值。

（3）双源条件下，矿井 1 的最晚充满时刻缩短约 $50.00\%$；矿井 2 新增 13 个可达端点和 14 条充满巷道。独立传播层使延迟水源不会被节点最早到达标签吞没。

（4）二次通知后，四名矿工改变出口。与继续沿 Q2 原路线相比，动态调整节省 $1.73$--$7.62\,\mathrm{min}$，且六条完整路径均通过连续水态复演。

守恒误差、手算小图、反例回归、官方模板合同和四因素灵敏度共同表明：该模型在题面信息范围内具有较好的正确性、稳定性和可复现性；其主要改进方向是引入实测水力参数，对平均分流规则进行标定。
